% Generated draft: requires mathematical and source review.
\documentclass[11pt]{article}
\usepackage[margin=25mm]{geometry}
\usepackage{fontspec}
\setmainfont{Latin Modern Roman}
\usepackage{amsmath,amssymb,mathtools,unicode-math}
\setmathfont{Latin Modern Math}
\usepackage{enumitem,needspace,xurl,hyperref}
\setlength{\parindent}{0pt}
\setlength{\parskip}{5pt}
\setlength{\emergencystretch}{3em}
\newcommand{\N}{\mathbb{N}}
\newcommand{\Z}{\mathbb{Z}}
\newcommand{\Q}{\mathbb{Q}}
\newcommand{\R}{\mathbb{R}}
\newcommand{\C}{\mathbb{C}}
\newcommand{\F}{\mathbb{F}}
\newcommand{\A}{\mathbb{A}}
\newcommand{\PP}{\mathbb{P}}
\newcommand{\E}{\mathbb{E}}
\newcommand{\eps}{\varepsilon}
\newcommand{\dd}{\,\mathrm{d}}
\begin{document}
\section*{UnsolvedMath ID: 30005553}
\textbf{Source problem code:} catalogue entry
\subsection{Definitions and mathematical statement}
Let \(0<\lambda<\tfrac12\). In the notation of the supplied record, \(R_t^\lambda\) denotes the planar Bochner--Riesz means, \(\operatorname{Sp}(p,q)\) denotes the sparse-bound property under consideration, and \(T_2(\lambda)\) denotes the necessary-condition trapezoid in the plane of reciprocal exponents. Write \(S_\lambda=P_2P_3\) for its designated boundary segment.

The source-relative question is whether, for every \(\lambda\in(0,\tfrac12)\) and every admissible exponent pair \((p,q)\),
\[
  (1/p,1/q)\in S_\lambda
  \quad\Longrightarrow\quad
  R_t^\lambda\text{ satisfies }\operatorname{Sp}(p,q).
\]
Here ``admissible'' refers to the exponent domain intended by the report; that domain is not specified in the supplied record.

This formulation retains unresolved source notation. The supplied material gives neither the multiplier formula and parameter range for \(R_t^\lambda\), nor the defining inequality for \(\operatorname{Sp}(p,q)\), nor the coordinates of \(P_2,P_3\) or defining inequalities of \(T_2(\lambda)\). It also does not specify whether the sparse bound is uniform in \(t\), whether a maximal operator is involved, or whether both endpoints of \(P_2P_3\) are included. These details must be recovered from the cited report before the question becomes a self-contained mathematical statement.

As background, the supplied report excerpt asserts that the sparse bound is obtained for reciprocal-exponent pairs in \(T_2(\lambda)\setminus P_2P_3\), and describes the behavior on \(P_2P_3\) as open. This is retained as a supplied source claim, not an independently verified theorem or current literature assessment.
\subsection{Short English statement}
For every \(0<\lambda<\tfrac12\), determine whether planar Bochner--Riesz means satisfy all sparse bounds corresponding to the designated boundary segment \(P_2P_3\). The supplied record omits the definitions needed to specify the segment, sparse inequality, exponent domain, and scale quantifiers precisely.
\subsection{Sources}
\begin{itemize}
\item Supplied local UnsolvedMath registry record, permanent ID \(30005553\), source problem code \texttt{catalogue entry}, titled ``Boundary Sparse Bounds for Planar Bochner--Riesz Means.'' The record belongs to the Oberwolfach Reports Open Problems collection and records the proposed year as 2023; no individual proposer is supplied.
\item Supplied report citation: \emph{Incidence Problems in Harmonic Analysis, Geometric Measure Theory, and Ergodic Theory} (2023), Oberwolfach Reports, pp.\ 1397--1452, DOI 10.4171/owr/2023/25. \url{https://doi.org/10.4171/owr/2023/25}.
\item Supplied literature source, labelled ``OWR 2023 25'': \url{https://publications.mfo.de/bitstream/handle/mfo/4055/OWR_2023_25.pdf?sequence=1}.
\item Literature assessment date: \texttt{2026-08-22}. The supplied dated curation labels the problem open, states that the report explicitly leaves the planar boundary-segment sparse bound open, and says that no later proof was found. This assessment is supplied dated curation and has not been independently verified here.
\item The registry labels the statement \texttt{corrected\_verified} and says it was checked against the official report and preserves its hypotheses, notation, and mathematical target. That verification is a supplied provenance assertion; the report has not been independently consulted here.
\end{itemize}
\end{document}
